\lecture{10}{25. September 2025}{Internal Loading and Stress}

\begin{problemwithimage}{f7_1}{7-1}
  The shaft is supported by a smooth thrust bearing at $A$ and a smooth journal bearing at $B$.
  Determine the resultant internal loadings acting on the cross section at $C$.
\end{problemwithimage}

\begin{derivation}
  We start by finding the resultant of the distributed loading as:
  \begin{equation*}
    W = w L = \qty{2}{\m} \cdot \qty{600}{\N\per\m}  = \qty{1200}{\N}
  .\end{equation*}
  Which acts vertically downwards at $C$.

  We will now take a global moment equilibrium about journal bearing $B$.
  \begin{align*}
    \sum M_B &= 0 \\
    \qty{1200}{\N} \cdot \qty{2.5}{\m} - \qty{900}{\N} \cdot \qty{1.5}{\m} - A_y \cdot \qty{4.5}{\m} &= 0 \\
    A_y &= \frac{\qty{1200}{\N} \cdot \qty{2.5}{\m} - \qty{900}{\N} \cdot \qty{1.5}{\m} }{\qty{4.5}{\m} } \\
    A_y &= \qty{366,66}{\N}
  .\end{align*}

  Now, taking a force equilibrium of the whole structure in the $y$ direction gives us our reaction $B_y$.
  \begin{align*}
    \sum F_y &= 0 \\
    A_y + B_y - \qty{1200}{\N} - \qty{900}{\N} &= 0 \\
    B_y &= \qty{2100}{\N} - \qty{366.66}{\N} \\
    B_y &= \qty{1733,34}{\N}
  .\end{align*}

  Now, using the method of sections, we take a cut at $C$ and only consider the left side ($AC$).
  Using a force equilibrium in the $y$ direction we get
  \begin{align*}
    \sum F_y &= 0 \\
    \qty{366.66}{\N} - \qty{600}{\N} + V_C &= 0 \\
    V_C &= \qty{600}{\N} - \qty{366.66}{\N} \\
    V_C &= \qty{233,34}{\N}
  .\end{align*}

  And a moment equilibrium gives us
  \begin{align*}
    \sum M_C &= 0 \\
    \qty{600}{\N} \cdot \qty{0.5}{\m} - A_{y}\cdot \qty{2}{\m} + M_C &= 0 \\
    M_C &= \qty{366.66}{\N} \cdot \qty{2}{\m} - \qty{600}{\N} \cdot \qty{0.5}{\m} \\
    M_C &= \qty{433,32}{\N.\m}\\
  .\end{align*}
  As we have no applied torque or axial load $T_C = 0$ and $N_C = 0$.
\end{derivation}

\finalresult{
\begin{aligned}
  V_C &= \qty{233,34}{\N}  \\
  M_C &= \qty{433,32}{\N.\m}  \\
  N_{C} &= \qty{0}{\N}  \\
  T_{C} &= \qty{0}{\N.\m} 
.\end{aligned}
}

\begin{problemwithimage}{f7_12}{F7-12}
  Determine the average normal stress in rod $AB$ if the load has a mass of \qty{50}{\kg}.
  The diameter of the rod $AB$ is \qty{8}{\mm}.
\end{problemwithimage}

\begin{derivation}
  We start by finding the force produced by the load.
  \begin{equation*}
    W = mg = \qty{50}{\kg} \cdot \qty{9.81}{\m\per\square\second}  = \qty{490,5}{\N}
  .\end{equation*}
  
  We now take a force equilibrium at pin $A$ to get
  \begin{align*}
    \sum F_y &= 0 \\
    T_{AC} \cdot \frac{3}{5} - W &= 0 \\
    T_{AC} &= W \cdot \frac{5}{3} \\
    T_{AC} &= \qty{490.5}{\N} \cdot \frac{5}{3} \\
    T_{AC} &= \qty{817.5}{\N}
  \end{align*}
  and
  \begin{align*}
    \sum F_x &= 0 \\
    \frac{4}{5} T_{AC} - T_{AB} &= 0 \\
    T_{AB} &= \frac{4}{5} T_{AC} \\
    T_{AB} &= \frac{4}{5} \cdot \qty{817.5}{\N} \\
    T_{AB} &= \qty{654}{\N}
  .\end{align*}
  I.e. $AB$ is in tension with $F_{AB} = \qty{654}{\N}$.
  
  The average normal stress is just the force divided by the area:
  \begin{equation*}
    \sigma_{AB} = \frac{F_{AB}}{A_{AB}} = \frac{\qty{654}{\N}}{\pi \cdot \left( \qty{4}{\mm}  \right)^2 }  = \qty{13}{\MPa}
  .\end{equation*}
\end{derivation}

\finalresult{\sigma_{AB} = \qty{13}{\MPa}}
